\section{Gaps, classes and the census}\label{sec:buys}

\subsection{A gap is a property of the class}

\begin{proposition}\label{prop:class}
Each rule of Part~I (\cref{lem:monotone} used in both directions, \cref{thm:split,thm:ruleA,thm:ruleB}) preserves
the quantum model set \emph{and} the classical model set. Hence, if $\G$ and $\G'$ are connected by these rules,
then $\G$ has a QC gap if and only if $\G'$ has one.
\end{proposition}
\begin{proof}
$\Q(\G)=\Q(\G')$ and $\C(\G)=\C(\G')$, so $\Q(\G)\not\subseteq\C(\G)$ is the same statement for both.
\end{proof}

This is the whole reason for classifying. It should be compared with what a \emph{classical} equivalence gives,
which is less.

\begin{lemma}[The upward rule]\label{lem:upward}
Let $\C(\G')=\C(\G)$ and let $\G'$ be structurally above $\G$ (more edges, larger sources). If $\G$ has a QC gap,
so has $\G'$.
\end{lemma}
\begin{proof}
$\Q(\G')\supseteq\Q(\G)\not\subseteq\C(\G)=\C(\G')$, the first step by \cref{lem:monotone}.
\end{proof}

Whether a gap also passes \emph{downward} inside a classical class is not known, and
\cref{thm:hlp-four,thm:sixnode} show why it is not automatic: within one classical class the quantum model set can
shrink strictly.

\subsection{The four-node census in numbers}

The repository enumerates all structures on three and on four visible nodes up to relabelling, with every source
quantum, and builds two graphs on them (\code{metagraph\_advanced.py}). In the \emph{quantum metagraph} the edges
are structural dominance and the rules of Part~I; in the \emph{classical metagraph} they are structural dominance,
the classical edge rule and strong face splitting. The strongly connected components are proven equivalence
classes, quantum and classical respectively. ``Proven'' is meant literally: two different classes may still be
equivalent by an argument that no rule captures.

\begin{center}
\small
\renewcommand{\arraystretch}{1.2}
\begin{tabular}{@{}lrr@{}}
\toprule
 & three nodes & four nodes\\
\midrule
structures up to relabelling & 46 & 2809\\
classical classes & 8 & 396\\
\quad of which contain a structure without sources & 5 & 20\\
quantum classes with the rules of Part~I & 8 & \textbf{446}\\
\quad for comparison, with strong face splitting (unsound) & 8 & 438\\
\quad for comparison, with weak face splitting only & 8 & 461\\
input classes of the census, and the structures in them & 3,\ 5 & 426,\ 996\\
\quad with a proven QC gap & 3,\ 5 & 380,\ 935\\
\quad open & 0 & 46,\ 61\\
\bottomrule
\end{tabular}
\end{center}

\emph{Inputs.} A class that contains a structure classically equivalent to one without sources cannot have a gap:
its classical model set is everything that satisfies the conditional independences read off the graph, and
quantum models satisfy those as well \cite{HLP}. The other $426$ quantum classes are the inputs of the census.

\emph{How a gap is proved.} A few structures have a gap by a direct argument; they are the \emph{seeds}: the
instrumental scenario, the Evans scenario, the triangle, and the Bell scenario with independent or with linked
settings (\code{known\_QC\_gaps.py}). Every other proof is a chain of \emph{reductions} that ends at a seed. A
reduction is a map from a structure $\G$ to a smaller or simpler structure $\G_0$ together with a proof that a gap
of $\G_0$ implies a gap of $\G$. Such a proof always has two halves, and the repository names them:
\begin{itemize}
\item \emph{the quantum direction}: a quantum model of $\G$ can produce every quantum distribution of $\G_0$, in a
  form that has the extra observable properties the step relies on (for instance, one node perfectly predicts
  another);
\item \emph{the classical direction}: if a distribution with those properties were classical in $\G$, the
  corresponding distribution would be classical in $\G_0$.
\end{itemize}
The reductions used are: fixing a node to a constant, marginalising a node, postselecting on a node, stitching two
nodes into one, making a quantum source classical (\emph{degradation}), and a family of steps due to Fritz in
which a node whose outcome another node predicts loses the parents that the predictor does not see. They are
described in \code{manuscript/piggybacks.md}. An equivalence is the degenerate reduction in which both directions
are identities.

\emph{What changed with the sound rule.} Restricting face splitting to \cref{thm:split} separates five classes
into $13$ (\cref{sec:open-four}). All $13$ have a gap, by fixing a node to a constant or by the upward rule, so the
proven classes go from $372$ to $380$ and the open classes are the same $46$ as before, with the same $61$
structures.

\subsection{Gaps that are reached only through equivalence}\label{sec:only-via}

How much of the census actually rests on \cref{prop:class}? To find out, the search was run on each of the $996$
structures of the input classes separately, so that no step may replace a structure by an equivalent one
(\code{scripts/only\_via\_equivalence.py}).

\begin{center}
\small
\renewcommand{\arraystretch}{1.2}
\begin{tabular}{@{}p{8.2cm}rr@{}}
\toprule
what the search may use, besides the reductions & structures reached & not reached\\
\midrule
all seeds, classes as in the census & 935 & \\
all seeds, no equivalence of any kind & 933 & \textbf{2}\\
all seeds, classical equivalence in the upward rule only & 934 & \textbf{1}\\
\bottomrule
\end{tabular}
\end{center}

The numbers are the same for the ten-stage search and for the deeper search that is now the default of the
repository (one Fritz step that may grow the structure up to six nodes). The two structures that no reduction
reaches lie in one class, the Bell scenario with linked settings, shown in \cref{fig:bell-class}.

\begin{figure}[H]
\centering
\begin{tikzpicture}[mdag,scale=0.82]
  % in every picture: settings 0 (top left) and 1 (top right), outcomes 3 (bottom left) and 2 (bottom right)
  \newcommand{\bellnodes}{\vn{x}{0,1.5}{$0$} \vn{y}{2,1.5}{$1$} \vn{a}{0,0}{$3$} \vn{b}{2,0}{$2$} \qs{q}{1,-0.7} \feed{q}{a,b}}
  % --- G' (top left)
  \begin{scope}[shift={(0,4.4)}]
    \bellnodes \qs{l}{1,2.25} \feed{l}{x,y} \de{x}{y} \de{x}{a} \de{y}{b}
    \node[lab,font=\footnotesize\bfseries] at (1,2.95) {$\mathsf{G}'$};
  \end{scope}
  % --- S (top centre)
  \begin{scope}[shift={(6.6,4.4)}]
    \bellnodes \de{x}{y} \de{x}{a} \de{y}{b}
    \node[lab,font=\footnotesize\bfseries] at (1,2.25) {$\mathsf{S}$ (a seed)};
  \end{scope}
  % --- G (top right)
  \begin{scope}[shift={(13.6,4.4)}]
    \bellnodes \qs{l}{-0.85,0.75} \feed{l}{x,a} \de{x}{y} \de{x}{a} \de{y}{b}
    \begin{scope}[on background layer]
      \node[draw=oired, line width=1.1pt, rounded corners=3pt, fit=(x)(y)(a)(b)(q)(l), inner sep=4pt] (frameG) {};
    \end{scope}
    \node[lab,font=\footnotesize\bfseries,text=oired,anchor=south] at (frameG.north) {$\mathsf{G}$};
  \end{scope}
  % --- G'' (bottom left)
  \begin{scope}[shift={(0,-1.6)}]
    \bellnodes \qs{l}{1,2.25} \feed{l}{x,y} \de{x}{a} \de{y}{b}
    \node[lab,font=\footnotesize\bfseries] at (1,-1.4) {$\mathsf{G}''$ (common cause of the settings)};
  \end{scope}
  % --- H (bottom right)
  \begin{scope}[shift={(13.6,-1.6)}]
    \bellnodes \qs{l}{-0.85,0.75} \feed{l}{x,a} \de{x}{y} \de{y}{b}
    \begin{scope}[on background layer]
      \node[draw=oired, line width=1.1pt, rounded corners=3pt, fit=(x)(y)(a)(b)(q)(l), inner sep=4pt] (frameH) {};
    \end{scope}
    \node[lab,font=\footnotesize\bfseries,text=oired,anchor=north] at (frameH.south) {$\mathsf{H}$};
  \end{scope}
  % --- equivalences, labelled by the rule that proves them
  \draw[<->, line width=0.9pt, oigreen!70!black] (2.8,5.15) -- node[lab,above,align=center] {\cref{thm:split}:\\ source on $0\to1$} (6.0,5.15);
  \draw[<->, line width=0.9pt, oigreen!70!black] (9.2,5.15) -- node[lab,above,align=center] {\cref{thm:split}:\\ source on $0\to3$} (12.1,5.15);
  \draw[<->, line width=0.9pt, oigreen!70!black] (1.0,3.2) -- node[lab,right,align=left] {Rule A:\\ edge $0\to1$} (1.0,1.2);
  \draw[<->, line width=0.9pt, oigreen!70!black] (14.6,3.05) -- node[lab,right,align=left] {Rule A:\\ edge $0\to3$} (14.6,0.75);
\end{tikzpicture}
\caption{The class of the Bell scenario with linked settings: settings $0$ and $1$, outcomes $3$ and $2$. Green
arrows are equivalences, labelled by the rule that proves them. In the repository's notation:
$\mathsf{S}=$ \str{0→1, 0→3, 1→2; Q\{2,3\}}; $\mathsf{G}'=\mathsf S$ with \str{Q\{0,1\}};
$\mathsf{G}''=$ \str{0→3, 1→2; Q\{0,1\}, Q\{2,3\}}; $\mathsf{G}=\mathsf S$ with \str{Q\{0,3\}};
$\mathsf{H}=$ \str{0→1, 1→2; Q\{0,3\}, Q\{2,3\}}. The two framed structures are reached by no reduction.}
\label{fig:bell-class}
\end{figure}

How each member gets its gap:
\begin{itemize}
\item $\mathsf S$ is a seed. Its gap is Bell's theorem with correlated settings: the source of the outcomes is
  independent of both settings, so any classical model is local for $P(x_3,x_2\mid x_0,x_1)$.
\item $\mathsf{G}'$ and $\mathsf{G}''$, in which the settings share a quantum common cause, follow from seeds by
  degradation: making that common cause classical gives the Bell scenario with a classical link between the
  settings, which is a seed by the same direct argument.
\item $\mathsf G$ follows from $\mathsf S$ by the upward rule (\cref{lem:upward}), which uses the \emph{classical}
  equivalence of the two; with no equivalence at all it is not reached.
\item $\mathsf H$ is reached in no other way than through its quantum equivalence with $\mathsf G$ (Rule A), and
  from there with $\mathsf S$ (\cref{thm:split}).
\end{itemize}

In $\mathsf H$ the outcome $3$ learns its setting $0$ only through a quantum source. No reduction turns a source
into an edge: reductions remove nodes, remove parents or make sources classical. Turning that source into an edge
is exactly what Rule A does, because the setting is premeasurable.

Two further facts put this class in perspective (same script).
\begin{itemize}
\item \emph{The class is not reachable from the other seeds at all.} If the three seeds with linked settings are
  left out, no member of the class is proven, with or without equivalence ($379$ classes and $930$ structures
  remain proven). No reduction leads from a Bell scenario with linked settings to one with independent settings,
  since removing the link is not a reduction. The class needs its own direct argument; equivalence then spreads
  it to all five members.
\item \emph{Elsewhere equivalence is a convenience.} In every other proven class, each member is reached by
  reductions on its own. What the classes buy there is economy and certainty: $426$ problems instead of $996$, one
  certificate per class, and the assurance that counting classes does not count the same problem twice.
\end{itemize}

\subsection{A classical source that is as good as a quantum one}

The seeds are often stated with some sources classical, for instance the Bell scenario in which each party
receives its setting through a classical common cause. Structures with some sources declared classical are called
\emph{hybrids}; in a model of a hybrid the states of the classical sources are diagonal. Pre-measurement tells when
the declaration makes no difference.

\begin{proposition}[\code{QmDAG.lift}]\label{prop:lift}
Let $F$ be a source of a hybrid all of whose members, except at most one, are premeasurable. Then declaring $F$
classical or quantum gives the same quantum model set.
\end{proposition}
\begin{proof}
Let $F$ be quantum and apply \cref{lem:premeasure} to each premeasurable member. Every member then holds the
outcomes $\ell$ of all of them, and what is left of the state is a state $\rho_m(\ell)$ of the one remaining member
$m$, labelled by $\ell$. That member can prepare $\rho_m(\ell)$ itself once it knows $\ell$. So $F$ only needs to
distribute $\ell$, which is classical.
\end{proof}

With this, the nine Bell seeds of the repository are two classes (independent settings, linked settings) and the
three instrumental seeds are one.

\subsection{What a new quantum-valid rule would buy}

Three statements are left open by this note.
\begin{enumerate}[leftmargin=*]
\item The edge rule when an other member of the origin's source has two sources and no visible parent (the open
  box at the end of \cref{sec:edge}).
\item The five four-node splits of \cref{sec:open-four}.
\item Whether a gap passes downward inside a classical class.
\end{enumerate}
None of them can help the $46$ open classes, and this can be read off the classification. Of the $396$ classical
classes on four nodes, $374$ contain a single quantum class and $22$ contain several ($2$ to $9$); all $22$ are
entirely proven. Every open class is the only quantum class of its classical class. So even the full classical
rules, were they valid, would not connect an open class to a proven one. A rule that helps must be new classically
as well.

What equivalence already does for the open classes is to shrink them: the $46$ classes hold $61$ structures, and
the largest holds $12$. A gap for one member is a gap for all.
